\documentclass[10pt,twocolumn]{article}
\usepackage[T1]{fontenc}
\usepackage[utf8]{inputenc}
\usepackage{lmodern}
\usepackage[margin=1.8cm,top=2.4cm,bottom=2cm,columnsep=0.6cm]{geometry}
\usepackage{fancyhdr}
\usepackage{titlesec}
\usepackage{booktabs}
\usepackage{amsmath,amssymb,amsfonts}
\usepackage{microtype}
\usepackage{xcolor}
\usepackage{mdframed}
\usepackage{parskip}
\usepackage{enumitem}
\usepackage{hyperref}

\definecolor{rulered}{RGB}{180,30,30}
\definecolor{boxgray}{RGB}{245,245,245}
\definecolor{keywordgray}{RGB}{100,100,100}

\pagestyle{fancy}
\fancyhf{}
\fancyhead[L]{\textsc{\small Claude Mag}}
\fancyhead[C]{\small\textit{Machine Learning Research Digest}}
\fancyhead[R]{\small 2026-10-26}
\fancyfoot[C]{\small\thepage}
\renewcommand{\headrulewidth}{0.8pt}

\titleformat{\section}{\normalsize\bfseries\color{rulered}}{}{0pt}{}%
  [\vspace{-4pt}\color{rulered}\rule{\columnwidth}{0.4pt}]
\titlespacing{\section}{0pt}{8pt}{4pt}

\hypersetup{colorlinks=true,urlcolor=rulered,linkcolor=black}

\begin{document}
\setlength{\parindent}{0pt}
\setlength{\parskip}{4pt}

{\small\color{keywordgray}\textbf{Keywords:}
  KV cache compression \textbullet{}
  mixture-of-experts \textbullet{}
  Fisher information \textbullet{}
  efficient LLM inference}\par\vspace{4pt}

{\LARGE\bfseries\raggedright
  MoEZip: Routing-Aware KV Cache Compression for Sparse
  Mixture-of-Experts LLMs}\par\vspace{4pt}

{\large\itshape\raggedright
  A Fisher-Information Routing Metric Cuts KV Memory by 10$\times$ While
  Preserving Expert Dispatch---Outperforming the Prior Art by 2.3--4.9$\times$}\par\vspace{6pt}

{\small\textbf{By} Minsang Kim, Seungjun Baek}\par\vspace{2pt}

{\small\color{keywordgray}arXiv:2610.30206 \textbullet{}
  NeurIPS 2026}
\par\vspace{2pt}\noindent{\color{rulered}\rule{\columnwidth}{0.6pt}}\vspace{6pt}

Sparse mixture-of-experts (MoE) language models are uniquely
vulnerable to KV cache eviction: removing context entries does not merely
reduce attention coverage---it changes which experts are selected at each
decode step, corrupting generation in ways that attention-only importance
scores cannot predict. MoEZip closes this gap with a
Fisher-information-based routing sensitivity metric that, combined
with instruction-token attention weight, identifies which KV entries
can be safely dropped. At a 10\% retention ratio, MoEZip
surpasses the prior state-of-the-art by $+24$ and $+20$ average
points on LongBench and RULER---a $2.3\times$ and $4.9\times$
multiplicative improvement---without any model fine-tuning.

\begin{mdframed}[backgroundcolor=boxgray,linecolor=rulered,linewidth=1pt,
    innerleftmargin=6pt,innerrightmargin=6pt,
    innertopmargin=6pt,innerbottommargin=6pt]
\textbf{\small AT A GLANCE}\par\vspace{4pt}
\begin{description}[leftmargin=0pt,labelsep=4pt,itemsep=2pt,parsep=0pt,
    font=\small\bfseries]
  \item[Problem:] KV eviction in sparse MoE LLMs alters which experts
    are selected at decode time; attention-only scores are blind to this
    routing perturbation.
  \item[Why it matters:] MoE models (Mixtral, DeepSeek-MoE, etc.) are
    the dominant efficiency paradigm; deploying them at long contexts
    requires KV compression, yet existing methods degrade generation quality.
  \item[Method:] Routing sensitivity via per-position Fisher information of the
    router distribution, combined with instruction-token attention scores, gives
    a dual-signal importance measure for each KV entry.
  \item[Result:] At 10\% KV retention, MoEZip scores $2.3\times$ higher
    on LongBench and $4.9\times$ higher on RULER than the prior SoTA (KVZip).
  \item[Evidence:] Evaluation on LongBench, RULER, and NIAH across both
    softmax-type and sigmoid-type MoE routing architectures.
\end{description}
\end{mdframed}

\section{The Big Picture}

Sparse MoE models reduce per-token floating-point operations by routing
each token to a small subset of expert feed-forward networks. The key
computational saving is that most parameters remain dormant on any given
token. However, memory at inference scales with context length through the
KV cache, and reducing it through eviction is essential for long-context
deployment.

The standard approach---evict entries with the lowest cumulative attention
weight---works well for dense transformers, where every parameter is active
for every token. It fails silently in MoE models because the router is a
function of the hidden state produced by the attention layer, and the
hidden state depends on which KV entries are present. Removing a
seemingly low-attention entry can shift the attention distribution enough
to change the router's top-$k$ expert selection, propagating an error that
compounds across all subsequent layers.

\section{Why Existing Approaches Fall Short}

\textbf{Attention-only scoring (H2O, SnapKV, PyramidKV):}
These methods score KV entries by the attention weight they accumulate across
query tokens. They are designed for dense transformers and have no term that
reflects the router's sensitivity to perturbations in the attention output.

\textbf{Magnitude-based pruning:}
Low-magnitude entries may reside in positions where the routing gate
function is highly sensitive; magnitude and routing sensitivity are not
correlated in general.

\textbf{KVZip:}
The prior SoTA for MoE KV compression, but its importance score does not
account for routing distribution shift caused by eviction, leaving a
systematic gap that MoEZip addresses.

\section{Core Mechanism}

MoEZip scores each KV entry $(i, l)$ (context position $i$, layer $l$) with
two signals:

\begin{enumerate}[itemsep=2pt,parsep=0pt]
  \item \textbf{Routing sensitivity $\mathcal{F}_l(i)$:} How much does
    evicting position $i$ shift the expert selection probabilities at layer $l$?
    Derived from the Fisher information of the routing distribution with respect
    to the attention output at position $i$.
  \item \textbf{Instruction attention $\alpha_l(i)$:} The accumulated attention
    weight that instruction tokens (the task-relevant prefix) assign to
    position $i$ at layer $l$.
\end{enumerate}

The final importance score is $\phi_l(i) = \alpha_l(i) + \lambda \cdot \mathcal{F}_l(i)$,
with $\lambda$ balancing the two signals. Entries with the lowest $\phi_l(i)$ are
evicted. No gradient computation or additional data is required.

\section{Technical Formulation}

Let $p_l(i) = \text{softmax}(W_r h_l(i))$ be the expert routing distribution at
layer $l$, conditioned on the hidden state $h_l(i)$ produced by attending over
context position $i$. The Fisher-information routing sensitivity is:
\[
\mathcal{F}_l(i) = \mathbb{E}_{e \sim p_l(i)}\!\left[\left(\nabla_{h_l(i)} \log p_l(i)[e]\right)^2\right]
\]
For softmax routing, this simplifies to a closed form involving the variance of
the routing logits:
\[
\mathcal{F}_l(i) = \text{Var}_{e \sim p_l(i)}\!\left[W_r^{(e)} h_l(i)\right]
\]
where $W_r^{(e)}$ is the $e$-th row of the routing weight matrix. An analogous
closed form is derived for sigmoid-type routing. The combined score is:
\[
\phi_l(i) = \alpha_l(i) + \lambda \cdot \mathcal{F}_l(i)
\]
Entries with the lowest $\phi_l(i)$ values are evicted to meet the
retention budget.

\section{Learning or Inference Procedure}

MoEZip is entirely training-free. At inference:

\begin{enumerate}[itemsep=2pt,parsep=0pt]
  \item During the prefill phase, compute $\mathcal{F}_l(i)$ and
    $\alpha_l(i)$ for each KV position $i$ and layer $l$.
  \item Sort positions by $\phi_l(i)$ in descending order and retain the
    top-$k$ (where $k$ is the retention budget).
  \item Proceed with autoregressive decoding over the compressed KV cache.
\end{enumerate}

No gradient is computed; the Fisher computation reuses quantities already
available from the attention forward pass. The method is plug-in compatible
with any MoE LLM.

\section{What the Guarantee Says}

The Fisher-information derivation provides a local linear approximation to the
change in the routing distribution caused by evicting position $i$. Specifically,
the first-order routing perturbation $\Delta p_l(i)$ satisfies:
\[
\|\Delta p_l(i)\|_2^2 \approx \mathcal{F}_l(i) \cdot \|\Delta h_l(i)\|_2^2
\]
where $\Delta h_l(i)$ is the change in the hidden state caused by removing
entry $i$. Retaining entries in decreasing order of $\mathcal{F}_l(i)$ therefore
minimizes the worst-case routing perturbation norm at the given budget.

\section{Experimental Findings}

\begin{itemize}[itemsep=2pt,parsep=0pt]
  \item \textbf{LongBench (long-document understanding):} At 10\% KV retention,
    MoEZip scores $2.3\times$ higher than KVZip on the average metric across
    six long-document subtasks.
  \item \textbf{RULER (structured reasoning):} MoEZip achieves $4.9\times$ higher
    scores than KVZip at 10\% retention; the gains are especially large on
    multi-hop retrieval subtasks where expert dispatch plays a critical role.
  \item \textbf{NIAH (needle-in-a-haystack):} MoEZip successfully retrieves
    needles at compression ratios where all attention-only baselines fail.
  \item \textbf{Both routing types:} Results are consistent across softmax-type
    (Mixtral-style) and sigmoid-type (DeepSeek-style) MoE architectures.
  \item \textbf{High retention convergence:} At $\geq$50\% retention, all methods
    converge toward full-cache performance, confirming the gains are specific to
    the high-compression regime.
\end{itemize}

\section{Ablations and Interpretation}

\textbf{Routing-only scoring ($\lambda=1$, $\alpha=0$):}
Underperforms on tasks where instruction-token attention is the dominant signal;
confirms both terms are needed.

\textbf{Attention-only scoring ($\lambda=0$):}
Matches existing dense-transformer baselines, confirming that the routing-aware
term alone drives the MoEZip gains.

\textbf{Sensitivity to $\lambda$:}
Performance is robust to $\lambda \in [0.5, 2]$; the default $\lambda = 1$
performs well across all benchmarks.

\textbf{Computational overhead:}
The Fisher computation adds $< 5\%$ latency to the prefill phase and is negligible
relative to the decode speedup from the compressed KV cache.

\section*{Reference}

Minsang Kim, Seungjun Baek.
\textbf{MoEZip: Routing-Aware KV Cache Compression for Sparse
Mixture-of-Experts LLMs}.
arXiv:2610.30206. NeurIPS 2026.
\url{https://arxiv.org/abs/2610.30206}

\end{document}
